\chapter{Заключение}
\label{ch:conclusion}

В ходе выполнения курсовой работы решена задача разработки веб-приложения Fun-Walk
для персонализированного планирования пешеходных маршрутов.

\textbf{Основные результаты:}
\begin{enumerate}
  \item Проведён анализ предметной области и сформулирована двухэтапная постановка
        задачи: дискретная оптимизация на графе и непрерывная маршрутизация OSRM.
  \item Разработана математическая модель, включающая шестимерное признаковое
        пространство POI, нормированное скалярное произведение с вектором
        предпочтений, штраф за отклонение от коридора, формулу Haversine,
        привлекательностно-взвешенный граф и алгоритм Dijkstra.
  \item Спроектирована и реализована клиент-серверная архитектура на NestJS и React
        с REST API, интерактивной картой Leaflet и модулем planner, выделенным
        для алгоритмической логики.
  \item Проведено тестирование, подтвердившее функциональность системы и
        чувствительность маршрутов к изменению пользовательских приоритетов.
\end{enumerate}

\textbf{Практическая значимость} работы состоит в демонстрации применения
классических методов дискретной математики и геодезии в современной веб-разработке.
Исходный код открыт и документирован, что позволяет использовать проект как
учебный стенд на курсах «Алгоритмы», «Программирование», «Web-технологии».

\textbf{Направления дальнейшего развития:}
\begin{itemize}
  \item загрузка POI из Overpass API / OpenStreetMap;
  \item интеграция live-данных AQI (WAQI, OpenAQ);
  \item персистентное хранение в PostgreSQL;
  \item замена sorted-array Dijkstra на binary heap; сравнение с A*;
  \item калибровка параметров $\alpha$, $\eta$ на размеченных пользовательских данных;
  \item развёртывание собственного OSRM-инстанса для production.
\end{itemize}

Таким образом, поставленная цель достигнута, все задачи выполнены. Разработанное
приложение Fun-Walk представляет собой работоспособный прототип, сочетающий
математическую строгость модели с практической реализацией на TypeScript.
